\documentclass{article}
\usepackage{booktabs}
\usepackage{multirow}
\usepackage{array}

\title{Advanced LaTeX Table Features}
\author{htex Library}
\date{2026}

\begin{document}

\maketitle

\section{Multicolumn Spanning}

The \texttt{multicolumn} command allows cells to span multiple columns:

\begin{tabular}{|l|l|l|l|}
\hline
\multicolumn{2}{|c|}{\textbf{Merged Header}} & \multicolumn{2}{|c|}{\textbf{Another Merged}} \\
\hline
A1 & A2 & B1 & B2 \\
C1 & C2 & D1 & D2 \\
\hline
\end{tabular}

\section{Multirow Spanning}

The \texttt{multirow} command allows cells to span multiple rows:

\begin{tabular}{|l|c|r|}
\hline
\multirow{2}{*}{Spans 2 Rows} & Data 1 & Value 1 \\
& Data 2 & Value 2 \\
\hline
Row 2 & Middle & Right \\
\hline
\end{tabular}

\section{Partial Horizontal Lines with Cline}

Use \texttt{cline} to draw horizontal lines across only specific columns:

\begin{tabular}{|l|c|r|r|}
\hline
Header 1 & Header 2 & Header 3 & Header 4 \\
\cline{2-4}
Data A & Data B & Data C & Data D \\
Data E & Data F & Data G & Data H \\
\cline{1-2}
End 1 & End 2 & & \\
\hline
\end{tabular}

\section{Booktabs Professional Styling}

The booktabs package provides three commands for professional-looking tables:

\begin{tabular}{lll}
\toprule
\textbf{Column 1} & \textbf{Column 2} & \textbf{Column 3} \\
\midrule
Data 1 & Data 2 & Data 3 \\
Data 4 & Data 5 & Data 6 \\
Data 7 & Data 8 & Data 9 \\
\bottomrule
\end{tabular}

\section{Complex Tables with Multiple Features}

Combining multirow, multicolumn, and booktabs:

\begin{tabular}{lccr}
\toprule
\multicolumn{2}{c}{\textbf{Category}} & \textbf{Value} & \textbf{Total} \\
\midrule
\multirow{2}{*}{Group A} & Item 1 & 10 & \multirow{2}{*}{30} \\
& Item 2 & 20 & \\
\midrule
\multirow{2}{*}{Group B} & Item 3 & 15 & \multirow{2}{*}{35} \\
& Item 4 & 20 & \\
\bottomrule
\end{tabular}

\section{Column Alignment Control}

Use \texttt{multicolumn} to override column alignment:

\begin{tabular}{|l|l|l|}
\hline
\multicolumn{1}{|c|}{Centered} & Left & Right \\
\hline
Normal Left & \multicolumn{1}{c|}{Override Center} & Normal Left \\
\hline
\end{tabular}

\section{Complex Mixed Table}

A realistic example combining all features:

\begin{tabular}{|l|c|c|c|r|}
\hline
\multicolumn{5}{|c|}{\textbf{Sales Report Q1-Q4}} \\
\hline
Region & Q1 & Q2 & Q3 & \multicolumn{1}{|c|}{Q4} \\
\cline{2-5}
North & 100 & 120 & 140 & 160 \\
\cline{2-4}
South & \multicolumn{2}{c|}{150} & 170 & 190 \\
East & 80 & 90 & \multirow{2}{*}{110} & 130 \\
West & 95 & 105 & & 125 \\
\hline
\end{tabular}

\end{document}
